\section{Exercícios integrados e projetos} \label{sec:ExercicioIntetradosProjetos}


\subsection*{Nível 1 -- Estruturas Numéricas Básicas}

\begin{enumerate}[label=\thesection.\arabic*]

\item Seja o vetor de medições de temperatura(°C) a seguir.
\[
t(^{\circ}\mathrm{C}) = \begin{bmatrix}
	37.454011884736246 \\
	95.071430640991622 \\
	73.199394181140505 \\
	59.865848419703660 \\
	15.601864044243651 \\
	15.599452033620265 \\
	5.808361216819947 \\
	86.617614577493512 \\
	60.111501174320878 \\
	70.807257779604555 \\
	2.058449429580245 \\
	96.990985216199434 \\
	83.244264080042171 \\
	21.233911067827616 \\
	18.182496720710063 \\
	18.340450985343381 \\
	30.424224295953771 \\
	52.475643163223786 \\
	43.194501864211574 \\
	29.122914019804192 \\
	61.185289472237947 \\
	13.949386065204184 \\
	29.214464853521815 \\
	36.636184329369172 \\
	45.606998421703594
\end{bmatrix}
\]
Este vetor contém 16 dígitos que representam registros térmicos de um forno industrial. 
Considerando que os dados deverão ser armazenados em precisão simples, realize as seguintes tarefas:
	\begin{itemize}
	\item Converta todos os elementos para o padrão IEEE 754 (precisão simples).
	\item Calcule o erro absoluto e relativo médio.
	\item Analise o impacto acumulado em somatórios.
	\end{itemize}
	%

	
\end{enumerate}
%\item Considere um vetor de posições $x(t)$ com 1000 pontos igualmente espaçados.
%
%\begin{itemize}
%\item Calcule a derivada usando:
%\begin{itemize}
%\item Diferença progressiva
%\item Diferença central
%\end{itemize}
%\item Compare erro e custo computacional.
%\end{itemize}
%
%\item Seja $P(t)$ um vetor com 1000 amostras de potência.
%
%\begin{itemize}
%\item Calcule a energia usando:
%\begin{itemize}
%\item Regra do Trapézio
%\item Regra de Simpson
%\end{itemize}
%\item Compare desempenho e precisão.
%\end{itemize}
%
%\item Dados experimentais $(x_i, y_i)$ com $n=500$:
%
%\begin{itemize}
%\item Ajuste um modelo linear via mínimos quadrados.
%\item Resolva o sistema normal associado.
%\item Avalie erro quadrático médio.
%\end{itemize}
%
%\end{enumerate}
%
%\subsection*{Nível 2 -- Sistemas e Modelagem Numérica}
%
%\begin{enumerate}[label=\thesection.\arabic*]
%
%\item Considere uma matriz $A \in \mathbb{R}^{100 \times 100}$.
%
%\begin{itemize}
%\item Resolva $Ax = b$ usando:
%\begin{itemize}
%\item Eliminação de Gauss
%\item Método de Jacobi
%\end{itemize}
%\item Compare número de operações ($\mathcal{O}(n^3)$ vs iterativo).
%\end{itemize}
%
%\item Seja um vetor $x \in \mathbb{R}^{500}$:
%
%\begin{itemize}
%\item Resolva $f(x_i) = e^{-x_i} - x_i$ para cada elemento.
%\item Utilize método de Newton vetorizado.
%\item Analise convergência global.
%\end{itemize}
%
%\item Dados $n=200$ pontos:
%
%\begin{itemize}
%\item Construa interpolação polinomial.
%\item Compare com ajuste por regressão.
%\item Discuta instabilidade para alto grau.
%\end{itemize}
%
%\item Considere $x(t)$ com 1000 pontos:
%
%\begin{itemize}
%\item Derive numericamente (velocidade e aceleração).
%\item Integre numericamente.
%\item Compare erro acumulado.
%\end{itemize}
%
%\end{enumerate}
%
%\section*{Nível 3 -- Problemas de Engenharia em Escala}
%
%\begin{enumerate}[label=\thesection.\arabic*]
%
%\item Avalie numericamente:
%
%\[
%f(x) = \frac{1 - \cos(x)}{x^2}
%\]
%
%para $x \in [10^{-8}, 1]$ com 1000 pontos.
%
%\begin{itemize}
%\item Analise erro de cancelamento.
%\item Proponha reformulação estável.
%\end{itemize}
%
%\item Ajuste não linear com $n=500$ pontos:
%
%\[
%y = ae^{bx}
%\]
%
%\begin{itemize}
%\item Linearize o problema.
%\item Resolva via mínimos quadrados.
%\item Compare com ajuste direto iterativo.
%\end{itemize}
%
%\item Resolva um sistema não linear com 50 incógnitas:
%
%\begin{itemize}
%\item Utilize método de Newton generalizado.
%\item Analise custo por iteração ($\mathcal{O}(n^3)$).
%\end{itemize}
%
%\item Projeto integrado:
%
%\begin{itemize}
%\item Gere dados experimentais com ruído ($n=1000$).
%\item Ajuste modelo polinomial.
%\item Resolva sistema linear associado.
%\item Integre numericamente.
%\item Derive numericamente.
%\item Analise propagação de erro IEEE 754.
%\end{itemize}

